% Figure: lifecycle carbon intensity of electricity by generation source,
% in grams of CO2-equivalent per kilowatt-hour, spanning the construction,
% fuel, operation, and decommissioning of each. Coal and gas dominate; the
% low-carbon sources cluster two orders of magnitude below, their small
% residue coming from manufacturing and construction rather than operation.
% Values are representative median figures (current as of 2024); the point is
% the two-order-of-magnitude gap, not the exact bar heights.
\begin{figure}[htbp]
\centering
\begin{tikzpicture}
\begin{axis}[
  width=13cm, height=7.5cm,
  xbar,
  xlabel={lifecycle carbon intensity ($\si{\gram}\,\text{CO}_2\text{eq}/\si{\kilo\watt\hour}$)},
  xmin=0, xmax=980,
  xtick={0,200,400,600,800},
  symbolic y coords={hydro,wind,nuclear,solar PV,geothermal,gas CCGT,coal},
  ytick=data,
  y dir=reverse,
  nodes near coords,
  nodes near coords style={font=\scriptsize},
  every node near coord/.append style={anchor=west},
  bar width=8pt,
  tick label style={font=\scriptsize},
  label style={font=\small},
  enlarge y limits=0.12,
]
\addplot[fill=engblue, draw=engblue] coordinates {
  (24,hydro) (11,wind) (12,nuclear) (43,solar PV) (38,geothermal) (490,gas CCGT) (820,coal)
};
\end{axis}
\end{tikzpicture}
\caption[Lifecycle carbon intensity of electricity by source]{Lifecycle
carbon intensity of electricity by generation source, in grams of
carbon-dioxide-equivalent per kilowatt-hour, counting construction, fuel,
operation, and decommissioning. Coal and combined-cycle gas stand two orders
of magnitude above the low-carbon sources, whose small residue comes from
manufacturing and construction (the concrete and steel of a dam, the silicon
of a panel, the towers of a wind farm) rather than from operation. The gap,
not the exact bar height, is the content: decarbonising the electricity
carrier means moving generation from the top two bars to the cluster at the
bottom. Representative median figures, current as of 2024.}
\label{fig:vol04ch14:lifecycle-carbon}
\end{figure}
